\section{Conditioning on survival and quasi-stationarity}
\label{sec:qsd}

Almost-sure extinction does not imply that every surviving realisation is small.
For a subcritical branching process the unconditional law concentrates at zero,
whereas the law conditional on \(Z_n>0\) may approach a stable non-degenerate shape.
That conditional limit is the object required here.

\begin{definition}[Quasi-stationary and Yaglom laws]
Let \(X_n\) be a Markov chain with absorbing state \(0\).  A probability distribution
\(\nu\) on the non-absorbing states is \emph{quasi-stationary} if
\begin{equation}
  \Prob_\nu(X_n=j\mid X_n\neq0)=\nu_j
  \qquad(n\geq0,\ j\neq0).
  \label{eq:qsd-definition}
\end{equation}
If, from a specified initial state,
\begin{equation}
  \Prob(X_n=j\mid X_n\neq0)\longrightarrow\nu_j,
  \label{eq:yaglom-definition}
\end{equation}
then \(\nu\) is the corresponding Yaglom limit.
\end{definition}

For a subcritical Galton--Watson process satisfying the usual non-degeneracy and
integrability conditions, the Yaglom limit exists; the original result is due to
Yaglom \cite{Yaglom1947}, with modern treatments in
\cite{Harris1963,AthreyaNey1972}.  The binary offspring law satisfies those
conditions.  We do not re-prove the general convergence theorem, but we derive the
survival scale and the first two limiting moments explicitly.  This distinction is
important: convergence of two moments is consistent with the Yaglom theorem, but
does not by itself prove convergence of the full conditional distribution.

\subsection{The late survival scale}
\label{sec:late-survival}

Fix \(0<p<1/2\) and put \(m=2p\in(0,1)\).  The survival recursion can be written
\begin{equation}
  S_{n+1}=mS_n\left(1-\frac{S_n}{2}\right).
  \label{eq:Sn-factorised}
\end{equation}
Since \(S_n\downarrow0\), the ratio \(S_{n+1}/S_n\) tends to \(m\).  The geometric
rate is therefore \(m^n\), but the ratio limit alone does not identify its
amplitude.  Define
\begin{equation}
  A(p)=\lim_{n\to\infty}\frac{S_n}{m^n}.
  \label{eq:A-definition}
\end{equation}
The product argument in \cref{sec:A-product} proves that this limit exists and lies
strictly between zero and one half.  Consequently,
\begin{equation}
  S_n\sim A(p)m^n
  \qquad(n\to\infty).
  \label{eq:late-survival-asymptotic}
\end{equation}
The relation is an asymptotic equivalence: the ratio of its two sides tends to one.
It is not a formula for the limit of \(S_n\), which is zero.

The constant \(A(p)\) retains the effect of the early nonlinear iterations in
\cref{eq:Sn-factorised}.  Once \(S_n\) is small, the factor \(1-S_n/2\) is close to
one and subsequent decay is nearly linear.  The accumulation of all earlier factors
is precisely what remains in \(A(p)\).

\subsection{The limiting conditional mean}
\label{sec:qsd-mean}

Because \(Z_n=0\) contributes nothing to its own expectation,
\begin{equation}
  \E[Z_n]
  =\Prob(Z_n>0)\E[Z_n\mid Z_n>0].
\end{equation}
Using \cref{eq:binary-pgf-mean,eq:Sn-definition} gives the exact identity
\begin{equation}
  \E[Z_n\mid Z_n>0]=\frac{m^n}{S_n}.
  \label{eq:conditional-mean-exact}
\end{equation}
The asymptotic \cref{eq:late-survival-asymptotic} therefore yields
\begin{equation}
  \E[Z_n\mid Z_n>0]\longrightarrow\frac{1}{A(p)}.
  \label{eq:qsd-mean}
\end{equation}
The geometric decay of the unconditional mean is exactly cancelled by that of the
survival probability.  The surviving population consequently has a finite limiting
mean even though the population as a whole becomes extinct almost surely.

\begin{figure}[htbp]
\centering
\begin{tikzpicture}
\begin{axis}[
  width=0.82\textwidth,height=7.0cm,
  axis lines=left,
  xlabel={generation \(n\)},
  ylabel={\(\E[Z_n\mid Z_n>0]\)},
  xmin=0,xmax=250,ymin=0,ymax=30,
  grid=major,grid style={dashed,gray!25},
  legend style={font=\small,at={(0.02,0.98)},anchor=north west},
  every axis plot/.append style={thick}]
  \addplot table[x=n,y=p030,col sep=comma]{figures/conditional_means.csv};
  \addlegendentry{\(p=0.30\)}
  \addplot table[x=n,y=p040,col sep=comma]{figures/conditional_means.csv};
  \addlegendentry{\(p=0.40\)}
  \addplot table[x=n,y=p045,col sep=comma]{figures/conditional_means.csv};
  \addlegendentry{\(p=0.45\)}
  \addplot table[x=n,y=p048,col sep=comma]{figures/conditional_means.csv};
  \addlegendentry{\(p=0.48\)}
  \addplot table[x=n,y=p049,col sep=comma]{figures/conditional_means.csv};
  \addlegendentry{\(p=0.49\)}
\end{axis}
\end{tikzpicture}
\caption{Conditional means computed from the exact recursion
\cref{eq:Sn-recursion,eq:conditional-mean-exact}.  Each curve approaches \(1/A(p)\).
As \(p\uparrow1/2\), both the limiting height and the time needed to approach it
increase.  The plotted data are generated by
\texttt{scripts/reproduce\_numerics.py}.}
\label{fig:conditional-means}
\end{figure}

\subsection{The limiting conditional variance}
\label{sec:qsd-variance}

Let \(M_n=\E[Z_n^2]\), and write \(\sigma^2=\Var(\xi)\).  Conditional on \(Z_n\),
the next generation is a sum of \(Z_n\) independent offspring counts, so
\begin{equation}
  \E[Z_{n+1}^2\mid Z_n]
  =\sigma^2Z_n+m^2Z_n^2.
\end{equation}
Taking expectations and using \cref{eq:gw-mean},
\begin{equation}
  M_{n+1}=m^2M_n+\sigma^2m^n,
  \qquad M_0=1.
  \label{eq:second-moment-recurrence}
\end{equation}
For the binary law, \(\sigma^2=4p(1-p)=m(2-m)\).  Solving
\cref{eq:second-moment-recurrence} for \(m\neq1\) gives
\begin{align}
  \E[Z_n^2]
  &=m^{2n}+\sigma^2m^{n-1}\frac{1-m^n}{1-m}
  \notag\\
  &=m^n\left(1+\frac{1}{1-m}\right)
    -\frac{m^{2n}}{1-m}.
  \label{eq:binary-second-moment}
\end{align}
The second line is specific to the binary variance.  It returns
\(\E[Z_1^2]=4p\), providing a quick sign check on the final term.

For \(m<1\), division by \(S_n\sim A(p)m^n\) yields
\begin{equation}
  \E[Z_n^2\mid Z_n>0]
  =\frac{\E[Z_n^2]}{S_n}
  \longrightarrow
  \frac{1}{A(p)}\left(1+\frac{1}{1-2p}\right).
  \label{eq:qsd-second-moment}
\end{equation}
Combining this with \cref{eq:qsd-mean}, the conditional variance satisfies
\begin{equation}
  \Var(Z_n\mid Z_n>0)
  \longrightarrow
  \frac{1}{A(p)}\left(1+\frac{1}{1-2p}\right)
  -\frac{1}{A(p)^2}.
  \label{eq:qsd-variance}
\end{equation}
These limits are the first two moments of the binary Yaglom law.  Both become large
near criticality, where the conditioning retains increasingly broad populations.

\subsection{A continuum approximation to the recursion}
\label{sec:riccati-approximation}

Although no approximation is needed for the preceding results, a continuum view of
the recursion helps to locate the near-critical time scale.  Put
\(\varepsilon=1-2p>0\).  The exact increment is
\begin{equation}
  S_{n+1}-S_n=-\varepsilon S_n-pS_n^2.
  \label{eq:Sn-exact-increment}
\end{equation}
Replacing the unit difference by a derivative gives the Riccati equation
\begin{equation}
  \frac{\dd S}{\dd n}=-\varepsilon S-pS^2.
  \label{eq:Sn-riccati}
\end{equation}
The substitution \(v=1/S\) gives \(v'=\varepsilon v+p\), and hence
\begin{equation}
  S(n)=\frac{\varepsilon}
  {C\varepsilon\e^{\varepsilon n}-p},
  \label{eq:Sn-riccati-solution}
\end{equation}
where \(C\) is fixed by matching to the early discrete dynamics.  At criticality,
\cref{eq:Sn-riccati} reduces to \(S'=-S^2/2\) and recovers the scale \(2/n\).
For \(\varepsilon>0\), it predicts a crossover on the scale
\(n\asymp1/\varepsilon\) followed by exponential decay.

The approximation cannot determine \(A(p)\): its integration constant is controlled
by precisely those early generations for which replacing a difference by a
derivative is least accurate.  We therefore use \cref{eq:Sn-riccati-solution} only
as a guide.  The exact product and series in \cref{sec:constant-A} recover the
missing amplitude and permit a rigorous near-critical analysis.

\FloatBarrier
